\documentclass{beamer}
\usepackage{multicol,hyperref}
\renewcommand{\arraystretch}{1.5}
\newcommand{\fttheory}[1]{\fcolorbox{blue}{blue}{\color{white}\textbf{#1}\hspace{\textwidth}}}
\newcommand{\ftexampl}[1]{\fcolorbox{orange}{orange}{\color{black}\textbf{#1}\hspace{\textwidth}}}
\newcommand{\ftdata}[1]{\fcolorbox{violet}{violet}{\color{white}\textbf{#1}\hspace{\textwidth}}}
\newcommand{\ftheader}[1]{\fcolorbox{lightgray}{lightgray}{\color{black}\textbf{#1}\hspace{\textwidth}}}
\begin{document}

\begin{frame}{\ftheader{The Efficient Market Hypothesis (EMH)}}
	\begin{itemize}
		\item The Efficient Market Hypothesis (EMH) states that prices of securities fully reflect available information about securities
		\item It highlights the difficulties associated with earning abnormal returns using either \textit{fundamental analysis} (predicting returns based on accounting or economic data) or \textit{technical analysis} (predicting returns based on the past history of returns). 
	\end{itemize}
\end{frame}

\begin{frame}{\fttheory{The Random Walk}}
	\begin{enumerate}
		\item A stock price follows a \textit{random walk} returns are unpredictable.
		\item Suppose that Apple traded for 266 USD at 12:00 and everyone knew that at the end of the day, it would trade for 270 USD (a risk-free return of 1.5\% in just one afternoon!)
		\item Everyone would start buying. The price would rise to 267 (a return of 1.13\%), to 268 (.75\%), and so on until at 12:01, when the price would be 270 (0\%). 
		\item Put differently, all information (such as the fact that the price will be 270 USD by the end of day). 
	\end{enumerate}
\end{frame}

\begin{frame}{\fttheory{The Random Walk}}
	\begin{enumerate}
		\item The random walk explains why the stock market does well on days when bad news comes out (e.g. the market expected that unemployment would rise by 1\%, traded stocks at price reflecting the expectation that unemployment would rise by 1\%, and then when unemployment rose by .95\%, the stock price rose).
	\end{enumerate}
\end{frame}

\begin{frame}{\fttheory{The Efficient Market Hypothesis (EMH)}}
	There are three forms of the hypothesis, depending on how much information it assumes investors have
	\begin{enumerate}
		\item\textbf{Weak Form:} the assertion that stockprices already reflect all information contained in the history of past trading
		\item\textbf{Semi-Strong Form:} the assertion that stockprices already reflect all publicly available information 
		\item\textbf{Strong Form:} the assertion that stockprices already reflect all relevant information including inside information
	\end{enumerate}
\end{frame}

\begin{frame}{\fttheory{Technical Analysis}}
	\begin{itemize}
		\item Technical analysis has its roots in Charles Dow's ``Dow theory'' (Dow, of the Dow Jones Industrial Average and the \text{Wall Street Journal}).
		\item Technical analysis claims that even if prices on average incorporate all publicly information, they deviate in the short-term in predictable ways.
		\item It's less about the long-term value of the security and more about its short-term supply and demand
		\item Technical analysts typically attribute these patterns to market psychology. 
	\end{itemize}
\end{frame}

\begin{frame}{\fttheory{Technical Analysis}}
	\begin{itemize}
		\item Common concepts in technical analysis include \textit{support} (a price below which a stock's price doesn't typically fall) and \textit{resistence} (a price above which a stock's price doesn't typically rise).  
		\item It follows that when a stock's price nears its support, one should buy (it can only go up) and when the price nears its resistence, one should sell (it can only go down).
		\item \textbf{Important Caveat:} one factor identified in research is momentum: specifically the return of a portfolio that goes long stocks that have down well and short stocks that have done poorly. While this certaintly has the flavor of technical analysis, the trading costs of such a strategy are high. 
		\item Technical analysis generates a lot of false positives (as evidence by the analysis of technicians in one study who made bold predictions about future prices. The prices were random). 
		\item Another indicator has to do when moving averages cross.
	\end{itemize}
\end{frame}

\begin{frame}{\ftexampl{Support and Resistence}}
 	Consider a stock ($S$) that pays a quarterly dividend of 1.2 USD per share and has a beta of 1.5. Suppose that the equity risk premium is 8\%. According to the CAPM, the discount rate should be
	\begin{equation}
		\text{E}[r_{S}]=r_{F}+\beta E[r_{M}-r_{F}]=1.5\times 8\%=12\%. 
	\end{equation}
	The dividend discount model predicts that the fair price is
	\begin{equation}
		P=\frac{D}{r_{S}}=\frac{1.2\text{ USD}}{12\%}=10\text{ USD}
	\end{equation}
	Suppose that the chart shows a support of 9 USD and a resistence of 11 USD and suppose that stock starts paying dividends of 1.8 USD per share. The fair price should now be 15 USD. Continued on next slide...
\end{frame}

\begin{frame}{\ftexampl{Support and Resistence}}
	If other traders follow the support/resistence logic, then as soon as the 
	
\end{frame}

\begin{frame}{\fttheory{Fundamental Analysis}}
	\begin{itemize}
		\item  
	\end{itemize}
\end{frame}

\begin{frame}{\ftexampl{Fundamental Analysis}}
	\item 
\end{frame}

\begin{frame}{\fttheory{Technical and Fundamental Analyses: A Summary}}
	\begin{itemize}
		\item Technical analysis can be easily automated and trading \textit{against} technical analysis is very easy to automate. 
		\item Technical analysis software need only the stock price data and a few milliseconds of computer time to generate trading orders.
		\item Example: if everyone sells at an arbitrary resistence level, then submit a limit buy order at that level; if everyone buys at an arbitrary support level, then submit a limit sell order at that level. 
		\item Fundamental analysis, in constrast, requires much more work, and the people who do that work are compensated for their effort. A casual trader may not be able to generate alpha, but a stock analyst who spends 80 hours a week studying a particular firm may very well be able to.
	\end{itemize}
\end{frame}

\end{document}
